\chapter{Endliche Folgen, Anfangsstücke und Rekursion}
\label{b21guide:ch:finite}

In diesem Kapitel ist $M$ eine beliebige Menge.
Alle Längen und Indizes sind natürliche Zahlen.

\section{Länge und Anfügen}

\begin{theorem}[Eindeutige Länge und leerer Randfall]
\label{b21guide:thm:length}
Ist derselbe Folgengraph sowohl eine Folge auf $[m]_0$ als auch
auf $[k]_0$, so ist $m=k$.
In jeder Menge $M$ gibt es genau eine leere Folge.
Eine Folge positiver Länge in $M$ existiert genau dann,
wenn $M\ne\varnothing$.
\end{theorem}

\begin{definition}[Anfangsstück und Anfügen]
Für $a\colon[m]_0\to M$ und $k\leq m$ ist
$a|_{[k]_0}=(a_0,\ldots,a_{k-1})$ das \emph{Anfangsstück
der Länge $k$}. Für $x\in M$ bezeichnet $a^\frown x$
die Folge auf $[m+1]_0$ mit
\[
 (a^\frown x)_i=
 \begin{cases}
 a_i,&i<m,\\
 x,&i=m.
 \end{cases}
\]
\end{definition}

\begin{theorem}[Eindeutiges Anfügen und Abspalten]
\label{b21guide:thm:append}
Zu einer Folge $a$ der Länge $m$ und $x\in M$ gibt es genau
eine Folge $b$ der Länge $m+1$, deren Anfangsstück der Länge
$m$ gleich $a$ und deren letztes Glied gleich $x$ ist:
$b=a^\frown x$.
Umgekehrt gilt für jede Folge $b$ der Länge $m+1$
\[
 b=\bigl(b|_{[m]_0}\bigr)^\frown b_m.
\]
Insbesondere sind Anfangsstück und letztes Glied eindeutig bestimmt.
\end{theorem}

\begin{theorem}[Wechsel zwischen Null- und Einsbasierung]
Eine Folge $(a_0,\ldots,a_{m-1})$ entspricht eindeutig
der auf $\{1,\ldots,m\}$ indizierten Folge
$(b_1,\ldots,b_m)$ mit $b_j=a_{j-1}$.
Zurück führt $a_i=b_{i+1}$. Für $m=0$ sind beide Folgen leer.
\end{theorem}

\section{Unendliche Folgen aus endlichen Anfangsstücken}

\begin{theorem}[Verträgliche Anfangsstücke bestimmen eine Folge]
\label{b21guide:thm:prefixes}
Für jedes $m\in\mathbb N$ sei eine Folge
$u^{(m)}\colon[m]_0\to M$ gegeben. Gelte
\[
 u^{(m+1)}|_{[m]_0}=u^{(m)}
 \qquad(m\in\mathbb N).
\]
Dann gibt es genau eine Folge $a\colon\mathbb N\to M$ mit
$a|_{[m]_0}=u^{(m)}$ für alle $m$.
Ihre Glieder sind $a_n=u^{(n+1)}_n$.
Umgekehrt bilden die endlichen Anfangsstücke jeder unendlichen
Folge eine solche verträgliche Familie.
\end{theorem}

\section{Rekursive Festlegung}

\begin{theorem}[Rekursion mit einem vorhergehenden Glied]
\label{b21guide:thm:recursion}
Seien $x\in M$ und $F\colon\mathbb N\times M\to M$ gegeben.
Dann gibt es genau eine Folge $a\colon\mathbb N\to M$ mit
\[
 a_0=x,\qquad a_{n+1}=F(n,a_n)\quad(n\in\mathbb N).
\]
Insbesondere bestimmen $x\in M$ und eine Schrittvorschrift
$T\colon M\to M$ genau eine Folge mit
$a_0=x$ und $a_{n+1}=T(a_n)$.
\end{theorem}

\begin{theorem}[Rekursion mit endlich vielen vorhergehenden Gliedern]
\label{b21guide:thm:higher-recursion}
Seien $r\geq1$, $x_0,\ldots,x_{r-1}\in M$ und
$F\colon\mathbb N\times M^r\to M$ gegeben;
$M^r$ bezeichnet die Menge der $r$-Tupel aus $M$.
Dann gibt es genau eine Folge $a\colon\mathbb N\to M$ mit
\[
 \begin{aligned}
 a_j&=x_j &&(0\leq j<r),\\
 a_{n+r}&=F(n,a_n,\ldots,a_{n+r-1}) &&(n\in\mathbb N).
 \end{aligned}
\]
\end{theorem}

\begin{theorem}[Rekursion mit dem gesamten bisherigen Anfangsstück]
Sei $M^{<\mathbb N}$ die Menge aller endlichen Folgen in $M$
und sei $G\colon M^{<\mathbb N}\to M$ gegeben.
Dann gibt es genau eine Folge $a\colon\mathbb N\to M$ mit
\[
 a_n=G\bigl(a|_{[n]_0}\bigr)\quad(n\in\mathbb N).
\]
Insbesondere ist $a_0=G(\varnothing)$.
\end{theorem}

Die Rekursionssätze benutzen die natürliche Indexmenge aus Band~10.
Aussagen über Existenz oder Eindeutigkeit beliebiger Funktionen
werden damit nicht erneut entwickelt.
